\subsection{PID namespace separato per la
verifica}\label{pid-namespace-separato-per-la-verifica}

Si vuole eseguire gli script di verifica direttamente
all\textquotesingle interno dell\textquotesingle ambiente di produzione
del team, isolando i processi utilizzando un PID namespace separato.

L\textquotesingle obiettivo e\textquotesingle{} quello di nascondere a
tool come \passthrough{\lstinline!pspy!}, \passthrough{\lstinline!ps!} e
\passthrough{\lstinline!top!} gli script eseguiti in fase di verifica,
evitando che il concorrente, dal suo ambiente possa vedere
l\textquotesingle exploit ed i suoi parametri di verifica.

Questa soluzione evita il costo di clonare il container, tutta via
e\textquotesingle{} inadatta al nostro scopo

\textbf{Il filesystem, resta completamente condiviso}:

\begin{itemize}
\tightlist
\item
  se un exploit agisce su un servizio web, questo potrebbe lasciare
  tracce dell\textquotesingle exploit nei log
\item
  se lo script di verifica o l\textquotesingle exploit crea cartelle
  temporanee, possono essere lette dal partecipante
\item
  se lo script carica il payload su una porta in cui il servizio
  vulnerabile e\textquotesingle{} esposto, allora con
  \passthrough{\lstinline!tcpdump!} posso ascoltare il traffico nel
  container.
\end{itemize}

Quando isoliamo il namespace bisogna anche:

\begin{itemize}
\tightlist
\item
  isolare il mount namespace, per isolare la cartella
  \passthrough{\lstinline!/proc!}
\end{itemize}

Tuttavia, tool come \passthrough{\lstinline!pspy!} tramite lo
\textbf{scanning degli eventi} \passthrough{\lstinline!inotify!} su
\passthrough{\lstinline!/proc!}, e lo \textbf{scan delle chiamate di
sistema}, puo\textquotesingle{} rilevare l\textquotesingle avvio di
nuovi processi.

Inoltre, l\textquotesingle exploit potrebbe essere di \textbf{tipo di
struttivo}, portando al \textbf{crash} di un servizio importante o
dell\textquotesingle intero sistema. Si vuole dunque evitare di
\textbf{interrompere il workflow} del team in fase di verifica.

Se e\textquotesingle{} abilitata la misurazione
dell\textquotesingle{}\textbf{SLA}, per la challenge, far crashare un
container vuol dire abbassare l\textquotesingle SLA e dunque penalizzare
il punteggio del team.

Per quanto riguarda la rete, un PID namespace separato non implica la
creazione di una nuova interfaccia di rete:

\begin{itemize}
\tightlist
\item
  Se lo script verificatore si trova sulla macchina vulnerabile, allora
  puo\textquotesingle{} mandare il payload utilizzando
  l\textquotesingle interfaccia di loopback, che e\textquotesingle{}
  accessibile anche ai partecipanti.
\item
  Il partecipante può eseguire \passthrough{\lstinline!tcpdump -i lo!} o
  ispezionare le connessioni attive con \passthrough{\lstinline!ss!} e
  \passthrough{\lstinline!netstat!}, catturando il traffico di verifica
  in chiaro.
\end{itemize}

Dunque, l\textquotesingle approccio basato su pid namespace
e\textquotesingle{} stato scartato, poiché non isola le altri componenti
del sistema.
